\chapter{Laurent Lafforgue}\label{ch2002laurentlafforgue:laurent-lafforgue-fields-medal-2002-a-comprehensive-report-on-his-work-and-impact}

\begin{center}\rule{0.5\linewidth}{0.5pt}\end{center}

\textbf{Prepared:} July 2026 \textbf{Subject:} Laurent Lafforgue (b. November 6, 1966, Antony, France) \textbf{Position:} Professor, Institut des Hautes \u00c9tudes Scientifiques (IH\u00c9S), Bures-sur-Yvette; CNRS researcher \textbf{Doctoral Advisor:} G\u00e9rard Laumon (Universit\u00e9 Paris-Sud) \textbf{Fields Medal Citation:} \emph{``For his proof of the Langlands correspondence for GL(r) over function fields.''}

\begin{center}\rule{0.5\linewidth}{0.5pt}\end{center}

\begin{prereqs} Algebraic geometry (stacks, moduli spaces, \u00e9tale cohomology), number theory (automorphic forms, L-functions, Galois representations), basic representation theory of reductive groups. \end{prereqs}

\begin{objectives} Explain the Langlands correspondence for GL(r) over function fields; describe the moduli space of shtukas and its compactification; understand how the Grothendieck--Lefschetz trace formula realizes the correspondence through cohomology; state the Ramanujan--Petersson conjecture as a consequence. \end{objectives}

\section{Problem History}

The Langlands program \cite{lafforgue2002}, initiated by Robert Langlands in a 1967 letter to Andr\u00e9 Weil, posits a vast web of conjectures connecting Galois representations to automorphic forms. The central \emph{reciprocity conjecture} asserts a bijection between irreducible $n$-dimensional $\ell$-adic representations of the absolute Galois group $\operatorname{Gal}(\overline{F}/F)$ and cuspidal automorphic representations of $\operatorname{GL}(n)$ over the adeles $\mathbb{A}_F$, preserving $L$-functions and $\epsilon$-factors. For $n=1$ this reduces to class field theory, proved in the early twentieth century.

For $\operatorname{GL}(2)$ over function fields, Vladimir Drinfeld proved the correspondence in 1974 using the theory of Drinfeld modules and elliptic sheaves, earning the 1990 Fields Medal. His construction was intrinsically rank-two: the moduli spaces of Drinfeld modules have dimension one, and the Eichler--Shimura relation controls their cohomology. Extending this to arbitrary rank $r \ge 3$ required fundamentally new geometric objects because no analogue of Drinfeld modules existed for higher rank.

The difficulty was twofold. First, one needed a suitable moduli space whose cohomology could carry both automorphic and Galois information simultaneously. Second, one needed a trace formula for that moduli space, generalizing the Arthur--Selberg trace formula from the adelic setting to a geometric setting. Both requirements demanded the development of new techniques in the theory of algebraic stacks and in the compactification of moduli spaces.

The function field case is more tractable than the number field case because function fields are fields of rational functions on algebraic curves over finite fields, providing a geometric framework in which cohomological tools---especially the Weil conjectures proved by Deligne---can be applied. Nonetheless, the rank-$r$ problem remained formidable for over two decades after Drinfeld's work.

A key difficulty specific to rank $r \ge 3$ is that the moduli of $r$-dimensional vector bundles on a curve has dimension $r^2(g-1)+1$, which grows quadratically in $r$. The geometry of these moduli spaces is far richer than the rank-two case, and the natural compactifications acquire complicated boundary structures that must be controlled. Furthermore, the Hecke operators---which encode the action of double cosets $K \backslash G(\mathbb{A}) / K$---become harder to analyze combinatorially as the rank increases, and the eigenvalue relations that encode the Langlands correspondence require entirely new algebraic methods.

For number fields, the situation is even more complex because the archimedean places introduce continuous spectra and the need for representation-theoretic methods at infinity. The function field case avoids these difficulties entirely: the ``places at infinity'' are just ordinary points on the curve, and the compactness of the curve provides a finite-dimensional framework. This is why the function field case was attacked first and why Lafforgue's methods, while powerful, do not directly transfer to the number field setting.

The $L$-function associated to an automorphic representation $\pi$ of $\operatorname{GL}(r)$ over a function field $F = \mathbb{F}_q(C)$ is an Euler product:
\[
L(s, \pi) = \prod_{x \in |C|} \frac{1}{\det(I - \operatorname{Frob}_x \, q^{-s} \mid V_\pi^I)^{-1}},
\]
where $V_\pi^I$ is the subspace of invariants under the Iwahori subgroup. The Langlands correspondence predicts that each local factor matches the characteristic polynomial of Frobenius acting on the corresponding Galois representation, and Lafforgue's proof establishes this matching globally.

\section{How Was It Solved}

Lafforgue's proof proceeds through five major steps. First, he constructs the moduli space $\operatorname{Cht}^r_{k}$ of \emph{chtoucas} (shtukas) of rank $r$ on a smooth projective curve $X$ over $\mathbb{F}_q$. A chtouca of rank $r$ with $k$ legs is a triple $(E, \operatorname{Fr}, \tau)$ consisting of a vector bundle $E$ of rank $r$ on $X$, an isomorphism $\operatorname{Fr}$ outside finitely many points, and a modification $\tau$ encoding the Frobenius structure. The moduli space is a Deligne--Mumford stack, proper over $X^k$, smooth of dimension $2r(k-1)$, equipped with a Hecke action at each point of $X$.

Second, Lafforgue constructs a compactification of $\operatorname{Cht}^r_k$ using toric geometry and the Hilbert--Mumford stability criterion. The compactified space is a smooth proper algebraic stack containing the original moduli space as an open dense substack. This compactification is essential for applying the trace formula in cohomology.

Third, Lafforgue proves a stabilization of the Arthur--Selberg trace formula for $\operatorname{GL}(r)$ over function fields, expressing the trace formula as a sum over stable conjugacy classes. The stabilization is the key technical step that allows the geometric side to be matched with the spectral side.

Fourth, Lafforgue applies the Grothendieck--Lefschetz trace formula to the compactified moduli space:
\[
\#\operatorname{Cht}^r_k(\mathbb{F}_{q^n}) = \sum_{i=0}^{2d} (-1)^i \operatorname{tr}(\operatorname{Fr}^n, H^i_c(\operatorname{Cht}^r_k)).
\]
The left side counts chtoucas over $\mathbb{F}_{q^n}$; the right side decomposes into contributions from automorphic representations. By comparing the two sides and using the stabilization, Lafforgue shows that the cohomology decomposes as:
\[
H^*(\operatorname{Cht}^r_k) = \bigoplus_{\pi} \pi^{\otimes k},
\]
where the sum runs over cuspidal automorphic representations $\pi$ of $\operatorname{GL}(r)$.

Fifth, the Frobenius morphism acts on this cohomology, yielding an $\ell$-adic representation of $\operatorname{Gal}(\overline{F}/F)$. The eigenvalues of Frobenius match the Hecke eigenvalues of the corresponding automorphic representation, establishing the Langlands correspondence. As a consequence, the eigenvalues of Frobenius are Weil numbers of absolute value $q^{1/2}$, proving the Ramanujan--Petersson conjecture for $\operatorname{GL}(r)$ over function fields.

The proof runs to several hundred pages, published primarily in \emph{Inventiones Mathematicae} (2002). The main paper, ``Chtoucas de Drinfeld et correspondance de Langlands,'' is over 300 pages long. A key intermediate result is the \emph{Fundamental Lemma} for $\operatorname{GL}(r)$, which relates orbital integrals on $\operatorname{GL}(r)$ to those on its endoscopic subgroups; this lemma was later proved in full generality by Ngô Bào Châu (Fields Medal 2010).

The \emph{Hecke algebra} $\mathcal{H}(\operatorname{GL}(r), K)$, where $K = \operatorname{GL}(r, \mathcal{O}_F)$ is the maximal compact subgroup, is the convolution algebra of $K$-bi-invariant functions on $\operatorname{GL}(r, \mathbb{A}_F)$. It is generated by the Hecke operators $T_x$ at each place $x$ of $F$, and the eigenvalues of these operators on an automorphic representation $\pi$ are the Satake parameters. Lafforgue's trace formula decomposes the Hecke algebra action on the cohomology of the moduli space of shtukas into simultaneous eigenspaces, each corresponding to an automorphic representation. The key identity is the \emph{character relation}: for each Hecke operator $T_x$, the trace of $T_x$ on $H^*(\operatorname{Cht}^r_k)$ equals the sum of the Satake eigenvalues weighted by the multiplicity of the corresponding automorphic representation.

\section{Impact of This Discovery}

Lafforgue's proof is the most significant result in the Langlands program since Drinfeld's work on $\operatorname{GL}(2)$. It provides a complete proof of the Langlands correspondence for $\operatorname{GL}(r)$ in the function field case and serves as a model for the geometric approach to the Langlands program. The techniques---particularly the construction and compactification of the moduli space of shtukas and the stabilization of the trace formula---can be adapted to other reductive groups. The proof confirmed a central prediction of the Langlands program and provided the definitive framework for the function field case, establishing the correspondence as a theorem rather than a conjecture for $\operatorname{GL}(r)$ over $\mathbb{F}_q(C)$.

The Ramanujan--Petersson conjecture, proved as a byproduct, has implications for the Sato--Tate conjecture and for the distribution of eigenvalues of Hecke operators. It establishes optimal bounds on the growth of automorphic forms, a central question in analytic number theory. Specifically, the conjecture predicts that for a cuspidal automorphic representation $\pi$ of $\operatorname{GL}(r)$ over a function field $F = \mathbb{F}_q(C)$, the Hecke eigenvalues $\lambda_\pi(x)$ at a place $x$ of $F$ satisfy $|\lambda_\pi(x)| \le d(r-1) \cdot q^{(r-1)/2}$, where $d(r-1)$ is the dimension of the $(r-1)$-th symmetric power representation of $\operatorname{GL}(r)$.

The \emph{Sato--Tate conjecture} for a modular form $f$ of weight $k$ on $\operatorname{GL}(2)$ over $\mathbb{Q}$ predicts that the normalized Fourier coefficients $a_p / p^{(k-1)/2}$ are equidistributed with respect to the Sato--Tate measure $\frac{2}{\pi} \sin^2 \theta \, d\theta$ on $[-1,1]$. Lafforgue's proof of the Ramanujan--Petersson conjecture for $\operatorname{GL}(r)$ over function fields implies the Sato--Tate conjecture in this setting, since the Ramanujan--Petersson bound $|\lambda_\pi(x)| \le d(r-1) \cdot q^{(r-1)/2}$ is the input needed for the equidistribution statement. The Sato--Tate conjecture for modular forms over $\mathbb{Q}$ was proved by Taylor and collaborators (2011), building on the techniques that Lafforgue developed for the function field case.

Beyond the specific theorem, Lafforgue's work introduced new techniques in algebraic geometry: the use of algebraic stacks for moduli problems, the compactification of moduli spaces using toric methods, and the application of the trace formula to \u00e9tale cohomology in the context of shtukas. These techniques have become standard tools in arithmetic geometry.

The proof also established the geometric Langlands program as a viable approach to the Langlands correspondence. The moduli space of shtukas is a geometric object whose cohomology realizes the classical Langlands correspondence, connecting the classical and geometric formulations of the program.

The \emph{Langlands program for function fields} is more tractable than for number fields because function fields have a geometric structure that allows the use of cohomological methods. The curve $X$ over $\mathbb{F}_q$ provides a geometric framework in which the Weil conjectures (proved by Deligne) can be applied, and the moduli spaces of shtukas have a rich geometric structure that encodes the Langlands correspondence. Lafforgue's proof exploits this geometric structure to establish the correspondence, and the techniques are specific to the function field setting. For number fields, the archimedean places introduce continuous spectra and the need for representation-theoretic methods at infinity, which are absent in the function field case.

Lafforgue's construction of the moduli space of shtukas also revealed a beautiful interplay between the geometry of the curve $X$ and the representation theory of $\operatorname{GL}(r)$. The Hecke operators, which act on the cohomology, are realized as correspondences between different strata of the moduli space, and their simultaneous eigenspaces correspond exactly to the automorphic representations appearing in the Langlands decomposition. This geometric picture provides a concrete realization of the abstract functoriality conjectures.

The \emph{geometric Satake correspondence}, due to Mirkovi\'{c}--Vilonen, establishes an equivalence between the category of perverse sheaves on the affine Grassmannian $\operatorname{Gr}_G$ and the category of representations of the Langlands dual group ${}^L G$. Lafforgue's work on shtukas is closely related to the geometric Satake correspondence: the moduli space of shtukas is a global analogue of the affine Grassmannian, and the Hecke action on the cohomology of the moduli space of shtukas is the global analogue of the convolution product on the affine Grassmannian. The geometric Satake correspondence provides the local ingredients for Lafforgue's global construction, and the interplay between the local and global geometries is a central theme of his work.

The \emph{automorphic side} of the Langlands correspondence for $\operatorname{GL}(r)$ consists of cuspidal automorphic representations, which are irreducible unitary representations of $\operatorname{GL}(r, \mathbb{A}_F)$ appearing in the discrete spectrum of $L^2(\operatorname{GL}(r, F) \backslash \operatorname{GL}(r, \mathbb{A}_F))$. Each such representation decomposes as a tensor product $\pi = \bigotimes_x \pi_x$ over all places $x$ of $F$, and the local components $\pi_x$ are irreducible admissible representations of $\operatorname{GL}(r, F_x)$. The Langlands correspondence matches each $\pi_x$ with a semisimple $r$-dimensional Weil--Deligne representation $\rho_x$ of $\operatorname{Gal}(\overline{F_x}/F_x)$, and Lafforgue's proof establishes this matching globally.

The \emph{Galois side} consists of $\ell$-adic representations $\rho: \operatorname{Gal}(\overline{F}/F) \to \operatorname{GL}(r, \overline{\mathbb{Q}_\ell})$, which are continuous homomorphisms from the absolute Galois group of $F$ to the general linear group over the $\ell$-adic numbers. The $\ell$-adic representations arise from the \u00e9tale cohomology of algebraic varieties over $F$, and they carry deep arithmetic information about the structure of $F$. The Langlands correspondence connects the arithmetic of Galois representations with the analysis of automorphic forms, and Lafforgue's proof provides the most complete realization of this connection in the function field case.

The \emph{Weil conjectures}, proved by Deligne (1974), assert that the \u00e9tale cohomology of a smooth projective variety over a finite field satisfies a Riemann hypothesis: the eigenvalues of Frobenius are Weil numbers of absolute value $q^{i/2}$ in degree $i$. Lafforgue's proof of the Ramanujan--Petersson conjecture for $\operatorname{GL}(r)$ can be viewed as an analogue of the Weil conjectures for the moduli space of shtukas: the eigenvalues of Frobenius on the cohomology of $\operatorname{Cht}^r_k$ are Weil numbers of absolute value $q^{1/2}$, reflecting the fact that the shtukas encode automorphic representations of $\operatorname{GL}(r)$. This analogy is precise: the Grothendieck--Lefschetz trace formula applied to $\operatorname{Cht}^r_k$ yields the same structure as the Weil conjectures, with the automorphic representations playing the role of the cohomology classes.

\section{Modern Views and Improvements}

Since 2002, Lafforgue has been working on extending his results from $\operatorname{GL}(r)$ to arbitrary reductive groups---the Langlands functoriality problem, which remains one of the central open problems in the field. The Hitchin fibration, which Lafforgue has studied extensively, provides an alternative geometric framework that may be more amenable to this generalization.

Vincent Lafforgue, Laurent's brother, proved the Langlands correspondence for $\operatorname{GL}(n)$ over function fields with additional structures, extending the ``automorphic to Galois'' direction for all $n$ simultaneously. The Lafforgue brothers are a rare example of two siblings both making fundamental contributions to the same area of mathematics.

For number fields, the Langlands correspondence is known for $\operatorname{GL}(1)$ (class field theory), $\operatorname{GL}(2)$ (Wiles, Taylor, Kisin), and for $\operatorname{GL}(n)$ over CM fields (Scholze, Shin). Lafforgue's proof for function fields remains a guiding model for how these results might be unified, though significant additional difficulties arise from archimedean places and the more complicated trace formula for number fields.

The \emph{Langlands dual group} ${}^L G$ of a reductive group $G$ is defined as the complex reductive group whose root system is the dual of the root system of $G$. For $G = \operatorname{GL}(r)$, the dual group is ${}^L G = \operatorname{GL}(r, \mathbb{C})$, and the Langlands correspondence relates representations of $G$ over $F$ to representations of ${}^L G$ over $\mathbb{C}$. The $L$-group ${}^L G \rtimes \operatorname{Gal}(\overline{F}/F)$ incorporates the Galois action, and the Langlands correspondence respects this semidirect product structure. Lafforgue's proof establishes the correspondence for $\operatorname{GL}(r)$ in the function field case, where the Galois group is the absolute Galois group of $\mathbb{F}_q(C)$ and the dual group is $\operatorname{GL}(r, \mathbb{C})$.

The \emph{local Langlands correspondence} for $\operatorname{GL}(n)$ over a local field $F$ assigns to each irreducible admissible representation $\pi$ of $\operatorname{GL}(n, F)$ an $n$-dimensional Weil--Deligne representation $\phi_\pi$ of the Weil group $W_F$. This local correspondence is a key ingredient in the global Langlands program, and Lafforgue's proof establishes it as a consequence of the global correspondence. The local correspondence preserves the local $L$-functions and $\epsilon$-factors, and it is compatible with parabolic induction and Jacquet functors.

The categorical formulation of the geometric Langlands program, developed by Beilinson and Drinfeld, reformulates the correspondence as an equivalence of categories between D-modules on $\operatorname{Bun}_G$ and quasi-coherent sheaves on $\operatorname{Loc}_{{}^L G}$. Lafforgue's work on shtukas connects the classical and geometric formulations, and recent progress by Gaitsgory and others continues to build on this foundation.

The \emph{shtuka} (or chtouca) is the geometric analogue of an automorphic form. A shtuka of rank $r$ on a curve $X$ over $\mathbb{F}_q$ is a triple $(E, \operatorname{Fr}, \tau)$ where $E$ is a rank-$r$ vector bundle on $X \times \operatorname{Spec}(\mathbb{F}_q[\operatorname{Fr}])$, $\operatorname{Fr}$ is an isomorphism outside finitely many points, and $\tau$ is a modification at the leg points. The moduli space of shtukas $\operatorname{Cht}^r_k$ with $k$ legs is a Deligne--Mumford stack of dimension $2r(k-1)$, and its points over $\mathbb{F}_{q^n}$ correspond to shtukas defined over $\mathbb{F}_{q^n}$. The Hecke operators act on $\operatorname{Cht}^r_k$ by modifying the bundle $E$ at a single point, and the simultaneous eigenspaces of all Hecke operators correspond to the automorphic representations appearing in the Langlands decomposition.

The \emph{compactification} of $\operatorname{Cht}^r_k$ that Lafforgue constructs adds boundary divisors corresponding to degenerations of the vector bundle. The boundary components are indexed by the combinatorics of lattice chains in $(F_\infty)^r$, where $F_\infty$ is the function field. The compactification is smooth and proper, and its boundary has normal crossings, which is essential for the application of the trace formula. The key technical innovation is the use of the Hilbert--Mumford criterion to construct the compactification as a GIT quotient, ensuring that the Hecke action extends to the boundary.

The \emph{stabilization of the trace formula} expresses the geometric side of the Arthur--Selberg trace formula as a sum of stable trace formulas over the endoscopic groups of $G$. For $G = \operatorname{GL}(r)$, the endoscopic groups are products $\operatorname{GL}(r_1) \times \cdots \times \operatorname{GL}(r_k)$ with $r_1 + \cdots + r_k = r$, and the stabilization involves comparing orbital integrals on $G$ with those on each endoscopic group. Lafforgue's stabilization is achieved by computing the trace formula on the compactified moduli space of shtukas, where the geometric structure of the boundary divisors provides the necessary identities between orbital integrals.

Recent work by Fargues and Scholze has established a ``geometrization'' of the local Langlands correspondence using the Fargues--Fontaine curve, suggesting a deeper geometric underpinning for the entire Langlands program. The relationship between Lafforgue's shtuka-based approach and the Fargues--Scholze perspective is an active area of research.

The \emph{Fargues--Fontaine curve} is a fundamental object in $p$-adic Hodge theory that provides a geometric bridge between $p$-adic Galois representations and vector bundles on a curve over a non-archimedean field. The local Langlands correspondence, which relates local representations of $\operatorname{GL}_n(F)$ to $n$-dimensional Weil--Deligne representations, can be reformulated as a statement about vector bundles on this curve. Fargues and Scholze showed that the completed cohomology of certain moduli spaces of shtukas realizes this correspondence geometrically, providing a new perspective on Lafforgue's classical results.

Lafforgue's approach to the Langlands program continues to inspire new developments. The theory of \emph{automorphic sheaves} on moduli spaces of shtukas provides a geometric framework for studying automorphic forms, and the \emph{Hecke eigensheaves} conjecture---which predicts that the Langlands correspondence can be realized as an equivalence of categories of sheaves---remains a central open problem in the geometric Langlands program.

\begin{center}\rule{0.5\linewidth}{0.5pt}\end{center}

\section*{Summary}
\addcontentsline{toc}{section}{Summary}

Lafforgue proved the Langlands correspondence for $\operatorname{GL}(r)$ over function fields by constructing the moduli space of shtukas, compactifying it, and applying the Grothendieck--Lefschetz trace formula to decompose its cohomology into automorphic representations. His work established the Ramanujan--Petersson conjecture in this setting and introduced geometric techniques that continue to shape the Langlands program.
